\documentclass[11pt]{article}

\usepackage{amsmath, amssymb}
\usepackage{geometry}
\usepackage{hyperref}
\usepackage{enumitem}

\geometry{margin=1in}

\title{Supplementary Note: DNA Methylation to Gene Programs and Gene Sets in dig-gene-set-extractors}
\author{}
\date{}

\begin{document}
\maketitle

\section*{Overview}
This note defines the methylation extractors implemented in \texttt{omics2geneset}:
\texttt{methylation\_cpg\_diff} and \texttt{methylation\_dmr\_regions}.
Inputs are differential CpG/probe or DMR tables, and outputs are gene programs plus GMT gene sets.

\section{Mental model}
DNA methylation changes are converted to a signed signal with configurable orientation.
If differential methylation is represented as $\Delta_p$ (higher means hypermethylation in condition A vs B),
define orientation factor $\eta \in \{-1,+1\}$ where:
\begin{itemize}[leftmargin=2em]
\item \texttt{delta\_orientation=activity\_oriented}: $\eta=-1$
\item \texttt{delta\_orientation=raw}: $\eta=+1$
\end{itemize}
Then:
\[
\tilde{x}_p = \eta \Delta_p.
\]
Under \texttt{activity\_oriented}, positive values indicate hypomethylation-associated permissive activity and negative values indicate hypermethylation-associated repressive activity.

With significance weighting:
\[
x_p = (\eta \Delta_p) \cdot \min\{-\log_{10}(p_p + \epsilon), C\},
\]
where $p_p$ is p-value/FDR, $\epsilon$ is a small floor, and $C$ is a cap.

\section{Linkage model axis}
Each CpG/DMR is treated as a genomic interval and linked to genes via $L_{pg}$:
\begin{itemize}[leftmargin=2em]
\item \texttt{promoter\_overlap}
\item \texttt{nearest\_tss}
\item \texttt{distance\_decay}
\item \texttt{external} (user-provided region-to-gene links)
\end{itemize}

\section{Program-family axis}
Implemented program families are:
\begin{itemize}[leftmargin=2em]
\item \texttt{promoter\_activity}: promoter-overlapping methylation signals
\item \texttt{distal\_activity}: non-promoter methylation signals
\item \texttt{linked\_activity}: all intervals linked to nearest genes (mostly QC/secondary)
\end{itemize}

Default preset is \texttt{connectable}, which emphasizes:
\begin{itemize}[leftmargin=2em]
\item \texttt{promoter\_activity} with \texttt{promoter\_overlap}
\item \texttt{distal\_activity} with \texttt{distance\_decay}
\end{itemize}

\section{Gene scoring and aggregation}
Given transformed feature weights $\alpha_p = \phi(x_p)$ and links $L_{pg}$,
feature contributions to gene $g$ are aggregated with default \texttt{weighted\_mean}:
\[
\text{num}_g = \sum_p \alpha_p L_{pg}, \qquad
\text{den}_g = \sum_p |L_{pg}|,
\]
\[
s_g = \frac{\text{num}_g}{\max(\text{den}_g, \epsilon)}.
\]

Alternative aggregations are \texttt{sum}, \texttt{mean}, and \texttt{max\_abs}.

\section{From gene scores to gene sets}
Selection is based on magnitude $|s_g|$ by default (for differential signal strength), while preserving signed $s_g$ in outputs.
Typical default:
\begin{itemize}[leftmargin=2em]
\item \texttt{select=top\_k}, \texttt{top\_k=200}
\item \texttt{normalize=within\_set\_l1}
\end{itemize}

Within-set normalization:
\[
w_g = \frac{|s_g|}{\sum_{h\in S} |s_h|}, \qquad g \in S.
\]

\section{GMT export and directionality}
Default GMT export uses signed splitting:
\begin{itemize}[leftmargin=2em]
\item \texttt{\_\_pos}: genes with positive score
\item \texttt{\_\_neg}: genes with negative score (ranked by magnitude)
\end{itemize}

Interpretation depends on \texttt{delta\_orientation}. Under \texttt{activity\_oriented}, \texttt{\_\_pos} corresponds to hypomethylation-associated activity increase and \texttt{\_\_neg} corresponds to hypermethylation-associated activity decrease. Under \texttt{raw}, signs follow the raw delta definition.

Default guardrails:
\begin{itemize}[leftmargin=2em]
\item \texttt{gmt\_min\_genes=100}
\item \texttt{gmt\_max\_genes=500}
\item small sets skipped unless \texttt{emit\_small\_gene\_sets=true}
\end{itemize}

\section{Practical caveats}
\begin{itemize}[leftmargin=2em]
\item Probe-level artifacts (cross-reactive probes, SNP overlap) should be filtered upstream when possible.
\item Coordinate conventions and genome build consistency are critical.
\item Distal methylation effects are context dependent; interpretation should be paired with orthogonal evidence where available.
\end{itemize}

\paragraph{Practical heuristics in implementation.}
The implementation emits explicit warnings when sex chromosomes are dropped, when a large fraction of rows cannot be mapped to GTF chromosome space, and when optional resources are requested but unavailable under skip policy. It also computes lightweight artifact diagnostics on emitted GMT sets (for example OR/KRT family dominance fractions) and warns when dominance exceeds a threshold. These diagnostics are advisory by default; exclusion filters are opt-in.

\end{document}
